\section{Εισαγωγή στα Μεγάλα Γλωσσικά Μοντέλα (\en{LLMs})}

\subsection{Εισαγωγή}

Τα Μεγάλα Γλωσσικά Μοντέλα (\en{LLMs}) έχουν αναδειχθεί σε μια τεχνολογία που αλλάζει ριζικά τον χώρο της τεχνητής νοημοσύνης, ιδιαίτερα στον τομέα της επεξεργασίας φυσικής γλώσσας (\en{NLP}). Τα μοντέλα αυτά μπορούν να κατανοούν, να παράγουν και να επεξεργάζονται την ανθρώπινη γλώσσα με εντυπωσιακή ακρίβεια. Η επίδρασή τους εκτείνεται σε πολλούς τομείς, όπως η παραγωγή κειμένου, η απάντηση ερωτήσεων, η παραγωγή κώδικα και πιο πρόσφατα, εφαρμογές κυβερνοασφάλειας, όπως η ανάλυση αρχείων καταγραφής και η ανίχνευση ανωμαλιών. Η γρήγορη εξέλιξη των \en{LLMs} έχει γίνει δυνατή χάρη στη διαθεσιμότητα συνόλων δεδομένων μεγάλης κλίμακας, στην αυξημένη υπολογιστική ισχύ και στην ανάπτυξη νέων αρχιτεκτονικών νευρωνικών δικτύων, κυρίως του μοντέλου \en{Transformer} \cite{vaswani2017attention}..

\subsection{Ορισμός των Μεγάλων Γλωσσικών Μοντέλων (\en{LLMs})}

Τα \en{LLMs} είναι μοντέλα βαθιάς μάθησης που εκπαιδεύονται σε τεράστιες ποσότητες κειμενικών δεδομένων, ώστε να μαθαίνουν τη στατιστική δομή της γλώσσας. Συνήθως αποτελούνται από δισεκατομμύρια ή ακόμη και τρισεκατομμύρια παραμέτρους, γεγονός που τους επιτρέπει να αποτυπώνουν σύνθετα γλωσσικά μοτίβα και σχέσεις που εξαρτώνται από τα συμφραζόμενα.

Τα \en{LLMs} λειτουργούν προβλέποντας την πιθανότητα μιας ακολουθίας λέξεων, με βάση το προηγούμενο συμφραζόμενο. Αυτή η ικανότητα τους επιτρέπει να εκτελούν ένα ευρύ φάσμα γλωσσικών εργασιών χωρίς να απαιτείται προγραμματισμός ειδικά για κάθε εργασία. Αντί γι’ αυτό, βασίζονται σε αναπαραστάσεις που έχουν μάθει μέσω προεκπαίδευσης σε ποικίλα σύνολα δεδομένων \cite{brown2020language}.

Τα βασικά χαρακτηριστικά των \en{LLMs} περιλαμβάνουν:
\begin{itemize}
\item Δεδομένα εκπαίδευσης μεγάλης κλίμακας
\item Υψηλή πολυπλοκότητα μοντέλου (εκατομμύρια έως δισεκατομμύρια παράμετροι)
\item Εφαρμοσιμότητα γενικού σκοπού σε διαφορετικές εργασίες
\item Κατανόηση γλώσσας με βάση τα συμφραζόμενα
\end{itemize}

Αυτά τα χαρακτηριστικά καθιστούν τα \en{LLMs} κατάλληλα για την ανάλυση μη δομημένων δεδομένων, όπως τα αρχεία καταγραφής δικτύου, όπου οι παραδοσιακές μέθοδοι συχνά δυσκολεύονται.

\subsection{Αρχιτεκτονική \en{Transformer}}

Η αρχιτεκτονική \en{Transformer}, η οποία παρουσιάστηκε από τους \en{Vaswani et al.} \cite{vaswani2017attention}, αποτελεί το βασικό δομικό στοιχείο των περισσότερων σύγχρονων \en{LLMs}. Σε αντίθεση με παλαιότερα μοντέλα ακολουθιών, όπως τα αναδρομικά νευρωνικά δίκτυα (\en{RNNs}), οι \en{Transformers} βασίζονται αποκλειστικά σε μηχανισμούς προσοχής για την επεξεργασία των δεδομένων εισόδου.

Σε γενικό επίπεδο, ο \en{Transformer} αποτελείται από μια δομή κωδικοποιητή-αποκωδικοποιητή, αν και πολλά \en{LLMs} χρησιμοποιούν μόνο ένα από αυτά τα δύο μέρη, ανάλογα με την εργασία. Η βασική καινοτομία του \en{Transformer} είναι ο μηχανισμός αυτοπροσοχής, ο οποίος επιτρέπει στο μοντέλο να σταθμίζει τη σημασία διαφορετικών λέξεων μέσα σε μια ακολουθία κατά τη δημιουργία αναπαραστάσεων.

Η αυτοπροσοχή υπολογίζει ταυτόχρονα τις σχέσεις μεταξύ όλων των \en{tokens} σε μια ακολουθία, επιτρέποντας στο μοντέλο να αποτυπώνει μακρινές εξαρτήσεις πιο αποτελεσματικά από τα παραδοσιακά μοντέλα. Αυτή η δυνατότητα παράλληλης επεξεργασίας βελτιώνει επίσης σημαντικά την αποδοτικότητα της εκπαίδευσης.

Τα κύρια συστατικά του \en{Transformer} περιλαμβάνουν:
\begin{itemize}
    \item Επίπεδα πολυκέφαλης αυτοπροσοχής (\en{multi-head self-attention layers})
    \item Νευρωνικά δίκτυα πρόσθιας τροφοδότησης (\en{feed-forward neural networks})
    \item Θεσιακή κωδικοποίηση (\en{positional encoding}) για τη διατήρηση της σειράς της ακολουθίας
    \item Κανονικοποίηση επιπέδων (\en{layer normalization}) και υπολειμματικές συνδέσεις (\en{residual connections})
\end{itemize}

Τα στοιχεία αυτά συνεργάζονται για να παράγουν ενσωματώσεις με βάση τα συμφραζόμενα, οι οποίες κωδικοποιούν τόσο τη σημασία των μεμονωμένων λέξεων όσο και τις σχέσεις τους μέσα στην ακολουθία.

\subsection{Προεκπαίδευση και Λεπτομερής Προσαρμογή}

Η επιτυχία των \en{LLMs} εξαρτάται σε μεγάλο βαθμό από τη μεθοδολογία εκπαίδευσής τους, η οποία συνήθως περιλαμβάνει δύο βασικά στάδια: την προεκπαίδευση και τη λεπτομερή προσαρμογή.

\subsubsection{Προεκπαίδευση}

Κατά την προεκπαίδευση, το μοντέλο εκπαιδεύεται σε μεγάλης κλίμακας, μη επισημασμένα κείμενα, χρησιμοποιώντας στόχους αυτοεποπτευόμενης μάθησης. Συνηθισμένες εργασίες προεκπαίδευσης περιλαμβάνουν τη γλωσσική μοντελοποίηση (πρόβλεψη της επόμενης λέξης σε μια ακολουθία) και τη μοντελοποίηση γλώσσας με μάσκα (πρόβλεψη λέξεων που λείπουν από μια πρόταση).

Αυτή η φάση επιτρέπει στο μοντέλο να μάθει γενική γλωσσική γνώση, όπως γραμματική, σύνταξη και σημασιολογικές σχέσεις. Επειδή η προεκπαίδευση δεν απαιτεί επισημασμένα δεδομένα, μπορεί να αξιοποιήσει τεράστιες ποσότητες δημόσια διαθέσιμου κειμένου \cite{devlin2019bert}.

\subsubsection{Λεπτομερής Προσαρμογή}

Μετά την προεκπαίδευση, το μοντέλο μπορεί να προσαρμοστεί λεπτομερώς σε ένα μικρότερο σύνολο δεδομένων, ειδικό για μια συγκεκριμένη εργασία. Η λεπτομερής προσαρμογή περιλαμβάνει τη ρύθμιση των παραμέτρων του μοντέλου, ώστε να βελτιστοποιηθεί η απόδοσή του για μια συγκεκριμένη εφαρμογή, όπως η ταξινόμηση, η απάντηση ερωτήσεων ή η ανίχνευση ανωμαλιών.

Η λεπτομερής προσαρμογή επιτρέπει στα \en{LLMs} να προσαρμόζονται σε συγκεκριμένους τομείς, συμπεριλαμβανομένης της κυβερνοασφάλειας, όπου μπορούν να εκπαιδευτούν ώστε να εντοπίζουν μοτίβα σε αρχεία καταγραφής δικτύου ή να ανιχνεύουν ύποπτες δραστηριότητες. Αυτή η διαδικασία εκπαίδευσης δύο σταδίων μειώνει σημαντικά την ανάγκη για μεγάλα επισημασμένα σύνολα δεδομένων σε μεταγενέστερες εργασίες.

\subsection{Παραδείγματα Μεγάλων Γλωσσικών Μοντέλων}

\subsubsection{\en{GPT} (\en{Generative Pre-trained Transformer})}

Η οικογένεια μοντέλων \en{Generative Pre-trained Transformer} (\en{GPT}) αποτελεί χαρακτηριστικό παράδειγμα \en{LLMs} που βασίζονται στην αρχιτεκτονική αποκωδικοποιητή \en{Transformer}. Τα μοντέλα \en{GPT} εκπαιδεύονται χρησιμοποιώντας έναν αυτοπαλίνδρομο στόχο γλωσσικής μοντελοποίησης, όπου το μοντέλο προβλέπει το επόμενο \en{token} σε μια ακολουθία με βάση το προηγούμενο συμφραζόμενο \cite{brown2020language}.

Τα μοντέλα \en{GPT} είναι ιδιαίτερα κατάλληλα για εργασίες παραγωγής κειμένου, όπως συστήματα διαλόγου, δημιουργία περιεχομένου και παραγωγή κώδικα. Η ικανότητά τους να παράγουν συνεκτικό και συναφές ως προς τα συμφραζόμενα κείμενο έχει οδηγήσει στην ευρεία υιοθέτησή τους τόσο στην ακαδημαϊκή κοινότητα όσο και στη βιομηχανία.

\subsubsection{\en{BERT} (\en{Bidirectional Encoder Representations from Transformers})}

Το \en{BERT} είναι ένα ακόμη σημαντικό \en{LLM}, βασισμένο στην αρχιτεκτονική κωδικοποιητή \en{Transformer}. Σε αντίθεση με το \en{GPT}, το οποίο επεξεργάζεται το κείμενο με μονοκατευθυντικό τρόπο, το \en{BERT} χρησιμοποιεί αμφίδρομο συμφραζόμενο, επιτρέποντάς του να λαμβάνει υπόψη τόσο τις προηγούμενες όσο και τις επόμενες λέξεις κατά τη δημιουργία αναπαραστάσεων \cite{devlin2019bert}.

Το \en{BERT} εκπαιδεύεται με τη χρήση μοντελοποίησης γλώσσας με μάσκα και εργασιών πρόβλεψης επόμενης πρότασης, γεγονός που το καθιστά ιδιαίτερα αποτελεσματικό στην κατανόηση της γλώσσας και όχι τόσο στην παραγωγή της. Έχει επιτύχει κορυφαίες επιδόσεις σε διάφορα σημεία αναφοράς \en{NLP}, όπως η ανάλυση συναισθήματος, η αναγνώριση ονοματισμένων οντοτήτων και η απάντηση ερωτήσεων.